\documentclass[12pt,spanish,letterpaper,oneside]{article}
\def\documenttitle {An\'alisis org\'anico y comparativo del Estado}
\def\documentsubtitle {El Estado frente al COVID-19: M\'exico, China y Estados Unidos}
\def\documenttitlehf {El Estado frente al COVID-19}
\def\hfwidthwrap {true}
\def\hfwidthtitle {0.48}
\def\hfwidthcourse {0.46}
\def\documentsubject {Actividad S7 - Cuadro comparativo}
\def\documentauthor {Martin Jonathan de la Cruz}
\def\coursename {Teor\'ia del Estado y Constituci\'on}
\def\coursecode {020}
\def\universityname {Universidad Abierta y a Distancia de M\'exico}
\def\universityfaculty {Licenciatura en Derecho}
\def\universitydepartment {Teor\'ia del Estado y Constituci\'on}
\def\universitydepartmentimage {departamentos/UnADM}
\def\universitydepartmentimagecfg {height=1.57cm}
\def\universitylocation {Ciudad de M\'exico}
\def\authortable {\begin{tabular}{ll}
\textbf{Alumno:} & Martin Jonathan de la Cruz \\
\textbf{Matr\'icula:} & ES2611202040 \\
\textbf{Figura docente:} & Ana Lilian S\'anchez Ch\'avez \\
\textbf{Semestre / Bloque:} & 2 / 2 \\
\textbf{Semana lectiva:} & Semana 7 \\
\textbf{Actividad oficial:} & Actividad S7: Cuadro comparativo (El Estado durante el COVID-19) \\
\textbf{Fecha de realizaci\'on:} & Noviembre de 2026 \\
\multicolumn{2}{l}{\universitylocation}
\end{tabular}}

\usepackage{iftex}
\ifLuaTeX
  \let\pdfpagewidth\pagewidth
  \let\pdfpageheight\pageheight
\fi

% Wrapper local para apuntar a la plantilla maestra de AulaTeX
\input{C:/Users/Sysx/Documents/AulaTeX-Academico/aulaTEX/base/Plantilla-Informe/template.tex}

\usepackage{helvet}
\renewcommand{\familydefault}{\sfdefault}
\setcitestyle{authoryear,open={(},close={)}}

\usepackage[most]{tcolorbox}
\usepackage{array,tabularx,booktabs,longtable,pdflscape}
\usepackage{tikz}
\usepackage{eso-pic}

\definecolor{unadmGreenDark}{HTML}{174A3A}
\definecolor{unadmGreen}{HTML}{2D6A4F}
\definecolor{tableHeader}{HTML}{E8EFE5}
\definecolor{rowAlt}{HTML}{F7FAF7}

\newcolumntype{L}[1]{>{\raggedright\arraybackslash}p{#1}}

\setlength{\emergencystretch}{3em}

\begin{document}
\AddToShipoutPictureBG*{%
  \begin{tikzpicture}[remember picture,overlay]
    \node[opacity=0.16,inner sep=0pt] at (current page.center)
      {\includegraphics[width=0.72\paperwidth]{img/departamentos/UnADM.pdf}};
  \end{tikzpicture}%
}

\templatePortrait
\templatePagecfg
\onehalfspacing

\begin{abstractd}
Las situaciones de emergencia sanitaria global representan el examen de esfuerzo m\'as severo para la teor\'ia del Estado contempor\'anea, pues someten a tensi\'on directa el principio de separaci\'on de poderes, los l\'imites del poder punitivo estatal y la vigencia irrenunciable de los derechos fundamentales. La crisis epidemiol\'ogica provocada por el virus SARS-CoV-2 (COVID-19) confront\'o a la comunidad internacional con dilemas estructurales sobre la distribuci\'on territorial del poder y los modelos de gesti\'on gubernamental. El presente estudio desarrolla un an\'alisis org\'anico y comparativo riguroso de tres paradigmas estatales contempor\'aneos frente a la pandemia: la Rep\'ublica representativa, democr\'atica y federal de los Estados Unidos Mexicanos; el Estado socialista unitario bajo centralismo democr\'atico de la Rep\'ublica Popular China; y la Rep\'ublica federal de frenos y contrapesos de los Estados Unidos de Am\'erica. Mediante una matriz anal\'itica integral de ocho dimensiones (concepto, forma de Estado, forma de gobierno, actuaci\'on del Poder Ejecutivo, actuaci\'on del Poder Legislativo, actuaci\'on del Poder Judicial, impacto hist\'orico y tutela efectiva de los derechos humanos), se eval\'uan las respuestas institucionales adoptadas y se articula una postura cr\'itica fundamentada en las concepciones del Estado de Derecho, concluyendo con lecciones esenciales para la subsistencia y perfeccionamiento del orden constitucional democr\'atico.\footnote{Declaraci\'on sobre el uso de herramientas de inteligencia artificial: En cumplimiento estricto con los Lineamientos para el uso \'etico y responsable de la inteligencia artificial de la Universidad Abierta y a Distancia de M\'exico (UnADM), se hace constar que se utiliz\'o un modelo de lenguaje artificial exclusivamente como apoyo instrumental en la correcci\'on de estilo y composici\'on editorial. La autor\'ia conceptual, el an\'alisis jur\'idico-dogm\'atico, la postura cr\'itica y las conclusiones son de exclusiva responsabilidad intelectual del estudiante.}
\end{abstractd}

\templateIndex
\templateFinalcfg

\section{Fundamentaci\'on dogm\'atica: Formas de Estado, formas de Gobierno y divisi\'on del poder}

El estudio cient\'ifico del Estado exige distinguir con precisi\'on conceptual entre las categor\'ias de \emph{forma de Estado} y \emph{forma de Gobierno}, dualidad te\'orica indispensable para comprender la anatom\'ia y la din\'amica funcional del poder pol\'itico \citep{andrade2003}.

La \textbf{forma de Estado} alude a la estructura total de la organizaci\'on pol\'itica en relaci\'on con sus elementos constitutivos b\'asicos: el pueblo, el territorio y la soberan\'ia. Modela la distribuci\'on espacial o territorial de las competencias soberanas, distinguiendo fundamentalmente entre el Estado unitario o centralizado ---donde impera un \'unico centro de decisi\'on pol\'itica y un orden jur\'idico uniforme para todo el territorio--- y el Estado federal ---donde la soberan\'ia se articula en una estructura compuesta y descentralizada, coordinando entidades federativas dotadas de autonom\'ia interna con una Constituci\'on Federal suprema \citep{kriele2022,tenaRamirez1999}.

Por su parte, la \textbf{forma de Gobierno} refiere al modo espec\'ifico en que se organizan, distribuyen y relacionan los \'organos supremos que ejercen las funciones del poder p\'ublico (Poder Ejecutivo, Poder Legislativo y Poder Judicial). Mientras la forma de Estado responde a la interrogante de \emph{c\'omo se reparte el poder en el espacio geogr\'afico y poblacional}, la forma de gobierno dilucida \emph{qui\'en ejerce el poder y bajo qu\'e sistema org\'anico de control rec\'iproco} \citep{bobbio1991,andrade2003}. Entre estas destacan el sistema presidencial, caracterizado por una estricta legitimaci\'on democr\'atica independiente de los poderes y la jefatura unificada de Estado y de Gobierno en el Presidente; el sistema parlamentario, basado en la colaboraci\'on de poderes y la confianza legislativa; y los sistemas de asamblea o centralismo unipartidista, donde el poder p\'ublico se unifica org\'anicamente bajo la direcci\'on ideol\'ogica del partido gobernante \citep{kelsen2020teoria}.

La crisis pand\'emica originada por el COVID-19 en 2020 activ\'o potestades extraordinarias en todos los Estados del mundo \citep{whoCovidPandemic2020}. En situaciones de cat\'astrofe biol\'ogica, el Estado constitucional enfrenta la tentaci\'on del hiperpresidencialismo y la expansi\'on desmedida de las medidas de polic\'ia sanitaria. Como advierte \citet{garciaRicci2011}, el principio de legalidad y el Estado de Derecho no admiten la claudicaci\'on de los derechos humanos bajo el pretexto de la emergencia. Todo acto de autoridad que restrinja prerrogativas fundamentales debe satisfacer las exigencias de reserva de ley, idoneidad, necesidad y proporcionalidad estricta \citep{carbonell2019,ferrajoli2001}. La comparaci\'on entre M\'exico, China y Estados Unidos permite contrastar c\'omo tres matrices pol\'iticas dispares respondieron al dilema entre seguridad sanitaria y libertad ciudadana.

\clearpage
\begin{landscape}
\thispagestyle{empty}
\section{Cuadro comparativo: El Estado durante el COVID-19}
\phantomsection\addcontentsline{toc}{section}{Cuadro comparativo: El Estado durante el COVID-19}

\begin{center}
\small
\begin{longtable}{p{4.2cm} p{7.3cm} p{7.3cm} p{7.3cm}}
\caption{Matriz comparativa org\'anica y funcional de actuaci\'on estatal frente a la pandemia de COVID-19}\\
\toprule
\rowcolor{tableHeader}
\textbf{Elemento anal\'itico} & \textbf{M\'EXICO} & \textbf{CHINA} & \textbf{ESTADOS UNIDOS (E.U.)} \\
\midrule
\endfirsthead
\toprule
\rowcolor{tableHeader}
\textbf{Elemento anal\'itico} & \textbf{M\'EXICO} & \textbf{CHINA} & \textbf{ESTADOS UNIDOS (E.U.)} \\
\midrule
\endhead
\bottomrule
\endfoot

\textbf{1. Concepto} & 
Rep\'ublica representativa, democr\'atica, laica y federal, soberana e independiente, regida por la Constituci\'on Pol\'itica de 1917, con primac\'ia garantista de los derechos humanos y divisi\'on tripartita del poder p\'ublico. & 
Estado unitario socialista de dictadura democr\'atica popular, dirigido por la clase trabajadora y fundamentado en la alianza obrero-campesina, bajo la direcci\'on monol\'itica del Partido Comunista Chino (PCCh). & 
Rep\'ublica federal constitucional y democr\'atica con r\'egimen presidencial y estricta separaci\'on de poderes, fundada en la soberan\'ia del pueblo y en la soberan\'ia residual de sus 50 Estados federados. \\
\midrule
\rowcolor{rowAlt}
\textbf{2. Forma de Estado} & 
\textbf{Estado Federal:} Compuesto por 32 entidades federativas libres y soberanas en su r\'egimen interior (Art. 40 CPEUM). En salubridad rige una concurrencia constitucional coordinada por el Consejo de Salubridad General (Art. 73, frac. XVI). & 
\textbf{Estado Unitario:} Estructura centralizada y vertical. Las provincias, regiones aut\'onomas y municipalidades carecen de soberan\'ia originaria y est\'an estrictamente subordinadas a las directrices del poder central de Beijing. & 
\textbf{Estado Federal:} Uni\'on indisoluble de 50 Estados con amplia autonom\'ia. Por mandato de la D\'ecima Enmienda, los Estados conservan el poder residual de polic\'ia (\emph{police power}) en materias de salud p\'ublica y orden interno. \\
\midrule
\textbf{3. Forma de Gobierno} & 
\textbf{Presidencialismo:} El titular del Ejecutivo Federal es Jefe de Estado y de Gobierno, electo por sufragio universal directo, conviviendo con un Congreso bicameral y un Poder Judicial independiente. & 
\textbf{Sistema de Asambleas Populares:} Centralismo democr\'atico unipartidista. La Asamblea Popular Nacional es formalmente el \'organo supremo, pero est\'a subordinada funcionalmente al Comit\'e Central del PCCh. & 
\textbf{Presidencialismo estricto:} R\'egimen de contrapesos rec\'iprocos (\emph{checks and balances}) entre el Presidente, un Congreso bicameral altamente fiscalizador y un Poder Judicial custodio de la Constituci\'on. \\
\midrule
\rowcolor{rowAlt}
\textbf{4. Poder Ejecutivo\newline(\textquestiondown C\'omo actu\'o?)} & 
$\bullet$ Declar\'o emergencia sanitaria por causa de fuerza mayor \citep{csgAcuerdoEmergencia2020}, evitando el estado de excepci\'on punitivo del Art. 29 CPEUM.\newline
$\bullet$ Coordinaci\'on v\'ia Secretar\'ia de Salud y Consejo de Salubridad General: Jornada de Sana Distancia, sem\'aforo epidemiol\'ogico y reconversi\'on hospitalaria masiva.\newline
$\bullet$ Confinamiento no coercitivo (sin toques de queda ni sanciones carcelarias a peatones); centralizaci\'on de compras y distribuci\'on gratuita universal de vacunas. & 
$\bullet$ Aplicaci\'on dogm\'atica de la pol\'itica ultrarrigorista de ``Cero COVID''.\newline
$\bullet$ Despliegue de confinamientos forzados inmediatos a ciudades enteras (Wuhan, Shangh\'ai), clausura de accesos, vigilancia geolocalizada masiva con c\'odigos QR de salud.\newline
$\bullet$ Construcci\'on r\'ecord de hospitales de campaña y aislamiento obligatorio de casos y contactos en centros gubernamentales; desarrollo masivo y exportaci\'on de vacunas. & 
$\bullet$ Declaraci\'on de emergencias nacionales bajo la Ley Stafford y Ley de Emergencias.\newline
$\bullet$ Lanzamiento de la \emph{Operaci\'on Warp Speed} con inversi\'on multimillonaria en vacunas de ARNm (Pfizer, Moderna).\newline
$\bullet$ Choque constante entre directrices federales y gobernadores estatales; emisi\'on de mandatos laborales federales de vacunaci\'on y mascarillas v\'ia agencias ejecutivas (OSHA). \\
\midrule
\textbf{5. Poder Legislativo\newline(\textquestiondown C\'omo actu\'o?)} & 
$\bullet$ El Congreso de la Uni\'on transit\'o a sesiones semipresenciales y virtuales.\newline
$\bullet$ Aprob\'o reasignaciones presupuestales de urgencia y la extinci\'on de fideicomisos p\'ublicos para canalizar fondos a la compra de vacunas y programas sociales.\newline
$\bullet$ Ejerci\'o una fiscalizaci\'on de intensidad moderada, manteniendo respaldo pol\'itico a la gesti\'on presidencial de la crisis sanitaria. & 
$\bullet$ La Asamblea Popular Nacional y su Comit\'e Permanente respaldaron un\'animemente todas las directrices del Bur\'o Pol\'itico del PCCh.\newline
$\bullet$ Aprobaci\'on expedita de legislaci\'on de bioseguridad y endurecimiento de sanciones por propagaci\'on de enfermedades o difusi\'on de noticias no autorizadas.\newline
$\bullet$ Inexistencia de oposici\'on parlamentaria o fiscalizaci\'on independiente. & 
$\bullet$ El Congreso (bicameral, con mayor\'ias pol\'iticas divididas) aprob\'o los mayores paquetes de est\'imulo fiscal de la historia moderna: \emph{CARES Act} (2.2 billones USD) y \emph{American Rescue Plan} (1.9 billones USD).\newline
$\bullet$ Estableci\'on de comit\'es de supervisi\'on rigurosa sobre el gasto de emergencia y comparecencias p\'ublicas constantes de autoridades sanitarias. \\
\midrule
\rowcolor{rowAlt}
\textbf{6. Poder Judicial\newline(\textquestiondown C\'omo actu\'o?)} & 
$\bullet$ El Poder Judicial de la Federaci\'on oper\'o mediante esquemas de guardia y aceler\'o la tramitaci\'on electr\'onica del juicio de amparo (FIREL) \citep{scjnAcuerdosCovid2020}.\newline
$\bullet$ Juzgados de Distrito otorgaron suspensiones de plano para garantizar atenci\'on m\'edica inmediata a infectados, entrega de equipos de protecci\'on a m\'edicos y enfermeras, y vacunaci\'on prioritaria a ni\~nas, ni\~nos y adolescentes. & 
$\bullet$ El Tribunal Popular Supremo y las fiscal\'ias subordinaron su actuaci\'on a los objetivos pol\'iticos de contenci\'on pand\'emica del r\'egimen.\newline
$\bullet$ Aplicaci\'on de justicia penal sumaria contra quienes violaron cuarentenas o divulgaron informaci\'on cr\'itica sobre el origen o manejo del brote.\newline
$\bullet$ Nula protecci\'on judicial aut\'onoma frente a detenciones o aislamientos arbitrarios. & 
$\bullet$ La Corte Suprema y los tribunales federales fungieron como \'arbitros constitucionales activos.\newline
$\bullet$ En \emph{NFIB v. OSHA} (2022) invalidaron el mandato federal de vacunaci\'on forzosa para empleados de grandes empresas por exceso competencial del Ejecutivo.\newline
$\bullet$ En \emph{Roman Catholic Diocese v. Cuomo} protegieron la libertad de culto contra restricciones desproporcionadas de aforos religiosos. \\
\midrule
\textbf{7. Impacto hist\'orico} & 
Contracci\'on del PIB del 8.5\,\% en 2020, severo impacto en el sector laboral informal y alta letalidad asociada a comorbilidades. Revel\'o la vulnerabilidad del sistema hospitalario p\'ublico, pero aceler\'o la digitalizaci\'on procesal judicial y afirm\'o la legitimidad social de no reprimir f\'isicamente a la poblaci\'on. & 
Demostr\'o una extraordinaria capacidad de movilizaci\'on infraestructural inicial que minimiz\'o contagios en 2020--2021; sin embargo, el encierro prolongado erosion\'o la econom\'ia, desat\'o hist\'oricas protestas de descontento popular (``folletos en blanco'') y culmin\'o en una apertura abrupta a finales de 2022. & 
M\'as de un mill\'on de decesos; exacerb\'o la polarizaci\'on ideol\'ogica y el debate sobre el rol del Estado federal; reconfigur\'o el r\'egimen laboral hacia el trabajo remoto global y consolid\'o el liderazgo estadounidense en la innovaci\'on biotecnol\'ogica farmac\'eutica. \\
\midrule
\rowcolor{rowAlt}
\textbf{8. \textquestiondown Se respetaron los Derechos Humanos?} & 
$\bullet$ \textbf{Aspectos garantizados:} Libertad de tr\'ansito preservada al evitar toques de queda y militarizaci\'on urbana punitiva; libertad de expresi\'on y cr\'itica plenamente vigentes.\newline
$\bullet$ \textbf{Vulneraciones/Tensiones:} El derecho a la salud (Art. 4.º CPEUM) sufri\'o por saturaci\'on hospitalaria y desabasto peri\'odico; la omisi\'on inicial de vacunar a menores motiv\'o cientos de juicios de amparo para hacer efectiva la protecci\'on. & 
$\bullet$ \textbf{Graves vulneraciones sistem\'aticas:} Supresi\'on total de la libertad de circulaci\'on mediante soldadura de accesos habitacionales; silenciamiento y criminalizaci\'on de m\'edicos y periodistas (caso Dr. Li Wenliang).\newline
$\bullet$ Invasiones masivas a la intimidad y privacidad v\'ia vigilancia biom\'etrica y confinamientos involuntarios en condiciones degradantes. & 
$\bullet$ \textbf{Garant\'ias activas:} Amplio escrutinio p\'ublico, libertad de prensa irreductible y tutela judicial efectiva de libertades individuales y de culto.\newline
$\bullet$ \textbf{Tensiones:} Fuerte conflicto entre autonom\'ia personal frente a mandatos sanitarios y el derecho colectivo a la vida; inequidad socioecon\'omica y racial en el acceso inicial a la atenci\'on hospitalaria. \\
\bottomrule
\end{longtable}
\end{center}

\end{landscape}
\clearpage

\section{An\'alisis cr\'itico y ponderaci\'on de los modelos estatales ante la emergencia sanitaria}

La valoraci\'on acad\'emica sobre qu\'e Estado asumi\'o de mejor manera la crisis del COVID-19 exige deslindarse de evaluaciones puramente tecnocr\'aticas o epidemiol\'ogicas cuantitativas. Desde la perspectiva de la \textbf{teor\'ia constitucional y el Estado de Derecho}, la excelencia de una respuesta estatal no puede medirse \'unicamente por la celeridad de sus medidas f\'isicas de aislamiento, sino por su capacidad para tutelar el bien jur\'idico de la salud p\'ublica \emph{preservando inc\'olume la dignidad humana y los l\'imites sustantivos del poder pol\'itico} \citep{garciaRicci2011,carbonell2019}.

\subsection{El colapso de la falacia autoritaria: La Rep\'ublica Popular China}

El modelo implementado por China ofreci\'o durante los primeros meses una aparente demostraci\'on de eficacia operativa: el confinamiento herm\'etico de Wuhan y el control biom\'etrico estricto redujeron dr\'asticamente la circulaci\'on del virus en su fase inicial. No obstante, una evaluaci\'on dogm\'atica integral evidencia que este resultado descans\'o sobre la destrucci\'on sistem\'atica de las libertades b\'asicas:
\begin{enumerate}
  \item \textbf{Opacidad y censura letal:} El silenciamiento penal y administrativo de los primeros m\'edicos que alertaron sobre el pat\'ogeno (como el doctor Li Wenliang) demor\'o la respuesta internacional e impidi\'o la prevenci\'on temprana, demostrando que la ausencia de libertad de expresi\'on es, en s\'i misma, un factor de riesgo epidemiol\'ogico.
  \item \textbf{Desproporcionalidad y trato degradante:} La pol\'itica de ``Cero COVID'' convirti\'o las medidas de salud en instrumentos de polic\'ia orwelliana. Encerrar a ciudadanos soldando las puertas de sus edificios y separar forzosamente a ni\~nas y ni\~nos de sus padres en campos de cuarentena constituye una violaci\'on flagrante de la dignidad y la integridad personal, incompatible con los est\'andares m\'inimos de derechos humanos \citep{scjnDignidad2016}.
  \item \textbf{Inflexibilidad estructural:} El r\'egimen unipartidista, carente de oposici\'on parlamentaria y de un Poder Judicial independiente, qued\'o atrapado en su propia ret\'orica de infalibilidad, postergando la transici\'on y desembocando en un colapso social y una reapertura desordenada a finales de 2022.
\end{enumerate}

\subsection{La fragmentaci\'on litigiosa y el choque de autonom\'ias: Los Estados Unidos}

En el extremo opuesto, los Estados Unidos exhibieron una extraordinaria potencia innovadora y cient\'ifica, patente en el financiamiento estatal r\'apido que permiti\'o a la industria farmac\'eutica desarrollar las vacunas de ARNm en tiempo r\'ecord mundial. Asimismo, su Congreso cumpli\'o una funci\'on ejemplar al autorizar est\'imulos econ\'omicos hist\'oricos que impidieron la ruina de millones de familias. 

Sin embargo, el sistema estatal estadounidense evidenci\'o graves deficiencias estructurales:
\begin{enumerate}
  \item \textbf{Hiperpolarizaci\'on y tribalismo partidista:} La gesti\'on de la salud p\'ublica fue instrumentalizada electoralmente. El escepticismo promovido desde sectores pol\'iticos respecto a las medidas b\'asicas de prevenci\'on fractur\'o la confianza institucional y deton\'o una resistencia virulenta a la vacunaci\'on, elevando la cifra de decesos a m\'as de un mill\'on de personas, la m\'as alta en t\'erminos absolutos del planeta.
  \item \textbf{Litigiosidad paralizante:} Aunque la intervenci\'on del Poder Judicial en fallos como \emph{NFIB v. OSHA} \citep{scotusNfib2022} reafirm\'o el principio de legalidad frente a la expansi\'on reglamentaria del Ejecutivo, la dispersi\'on judicial y los amparos cruzados generaron certidumbre fragmentada en la aplicaci\'on de directrices sanitarias indispensables.
\end{enumerate}

\subsection{El garantismo social y la preservaci\'on del orden constitucional: M\'exico}

En este marco comparativo, la respuesta del Estado mexicano, pese a enfrentar grav\'isimas restricciones presupuestales y una fragilidad hist\'orica en su infraestructura hospitalaria, presenta m\'eritos jur\'idicos e institucionales de primer orden que fundamentan una \textbf{valoraci\'on positiva superior en t\'erminos de legitimidad democr\'atica y respeto a las libertades civiles}:
\begin{enumerate}
  \item \textbf{Rechazo al estado de sitio policial:} La decisi\'on gubernamental de clasificar la contingencia como ``emergencia sanitaria por causa de fuerza mayor'' \citep{csgAcuerdoEmergencia2020} en lugar de solicitar la suspensi\'on formal de garant\'ias del art\'iculo 29 constitucional fue sumamente acertada. En una sociedad con m\'as del 50\,\% de su poblaci\'on activa en el sector informal, imponer toques de queda militarizados o sanciones penales a quienes sal\'ian a buscar el sustento diario habr\'ia generado violencia social incalculable y criminalizaci\'on de la pobreza. El confinamiento voluntario y solidario apel\'o a la corresponsabilidad ciudadana sin quebrar el Estado de Derecho.
  \item \textbf{Centralizaci\'on y universalidad de la protecci\'on biol\'ogica:} El Estado asumi\'o la adquisici\'on y distribuci\'on monop\'olica de vacunas de manera estrictamente p\'ublica y gratuita, evitando que la capacidad econ\'omica individual determinara el acceso a la inmunizaci\'on, materializando el principio de igualdad sustantiva.
  \item \textbf{El Juicio de Amparo como garant\'ia viva:} Lejos de suspenderse o anularse, el Poder Judicial de la Federaci\'on oper\'o como un contrapeso garantista efectivo. Frente a omisiones de las autoridades administrativas de salud ---como la negativa inicial a vacunar a la poblaci\'on menor de edad o la falta de insumos de protecci\'on a m\'edicos residentes---, los juzgados de distrito y los tribunales colegiados concedieron suspensiones protectoras que corrigieron la actuaci\'on del Ejecutivo federal, demostrando la operatividad real de la divisi\'on de poderes en M\'exico.
\end{enumerate}

\section{Conclusi\'on respecto a la decisi\'on estatal y lecciones para la Teor\'ia del Estado}

El an\'alisis org\'anico desarrollado a lo largo de este estudio permite formular tres conclusiones sustantivas sobre la teor\'ia del Estado frente a fen\'omenos de cat\'astrofe:

\begin{enumerate}
  \item \textbf{La indisolubilidad entre Estado de Derecho y pol\'itica sanitaria:} Los fines leg\'itimos de protecci\'on a la vida colectiva jam\'as pueden justificar la suspensi\'on arbitraria de los derechos fundamentales ni la entronizaci\'on de mecanismos tir\'anicos de polic\'ia. El r\'egimen que aniquila la libertad para ``salvar'' biol\'ogicamente a sus miembros termina destruyendo la raz\'on misma del pacto social.
  \item \textbf{El imperativo de la suficiencia material del Estado prestacional:} La pandemia evidenci\'o que la consagraci\'on formal de derechos sociales en la Constituci\'on (como el derecho a la salud del art\'iculo 4.º) resulta inoperante si el Estado no cuenta con capacidad presupuestal, reservas estrat\'egicas e instituciones hospitalarias s\'olidas. La soberan\'ia moderna no s\'olo es coercitiva, sino prestacional y cient\'ifica.
  \item \textbf{La vigencia salvaguarda de la divisi\'on de poderes:} Frente a la natural concentraci\'on de atribuciones que reclama el Poder Ejecutivo en coyunturas de emergencia, la existencia de un Poder Judicial t\'ecnico, independiente y accesible mediante v\'ias expeditas de control constitucional (como el juicio de amparo) constituye la \'ultima l\'inea de defensa de la persona frente a la desmesura estatal.
\end{enumerate}

\clearpage
\bibliography{teoria-del-estado-y-constitucion}
\bibliographystyle{apalike}

\end{document}
